\documentclass[11pt]{article}

\usepackage[a4paper,margin=1in]{geometry}
\usepackage{amsmath,amssymb,amsthm,mathtools}
\usepackage[hidelinks]{hyperref}
\hypersetup{
  pdfauthor={Bangyen Pham},
  pdftitle={Coefficient Mass of Polynomial Multiples in Sectors},
  pdfsubject={Lower bounds for the coefficient order statistics and logarithmic coefficient mass of polynomial multiples with prescribed roots in a sector},
  pdfkeywords={polynomial multiples, coefficient mass, complex roots, sectors, argument principle, Gaussian integers},
  pdflang={en-US}
}

\newtheorem{theorem}{Theorem}[section]
\newtheorem{lemma}[theorem]{Lemma}
\newtheorem{corollary}[theorem]{Corollary}
\newtheorem{proposition}[theorem]{Proposition}
\newcommand{\R}{\mathbb R}
\newcommand{\N}{\mathbb N}

\title{Coefficient Mass of Polynomial Multiples in Sectors}
\author{\shortstack{Bangyen Pham\\
\small Independent Researcher\\
\small\href{mailto:bangyenp@gmail.com}{\texttt{bangyenp@gmail.com}}\\
\small\href{https://orcid.org/0009-0006-9928-2088}{ORCID: 0009-0006-9928-2088}}}
\date{September 2026}

\begin{document}
\maketitle

\begin{abstract}
Let $b_k(F)$ be the $k$-th largest nonleading coefficient magnitude of
$F=\sum_{j\le D}f_jx^j$ and $\Lambda(F)=\sum_j\log^+|f_j|$.  For $K$
prescribed roots of modulus at least $R$ in a sector of half-width $\delta$,
we bound $\Lambda$ and the $b_k$ of multiples with $|f_D|\ge1$.  The
argument principle along a ray next to the sector gives many large real
zeros of a real polynomial dominated coefficientwise by the multiple, to
which the real-root bounds of a companion paper apply.  In thin sectors and
small degree this gives rows and mass of order $K^2\log R$.  Counting the
crossings against the equidistributed share $\delta/\pi$ of the large
zeros, multiples of degree $O(K)$ pay $(\frac\pi\delta-o(1))K\log R$ at
fixed $\delta<\pi/2$ and $(\frac12-o(1))K^2\log R$ when $\delta K\to0$, as
$K,R\to\infty$; both constants are sharp, and $x^N-\rho^N$ attains the
first.  At distinct Gaussian-integer roots in the upper half-plane, degree
charging removes the degree restriction: integer multiples of every degree
pay $(\frac\pi\delta-o(1))K\log\max(R,\sqrt K)$ as $K\to\infty$.  Two integer imitations of
$x^N-\rho^N$ are limited, and a small second row confines the roots to one
level set of a polynomial with small nonleading coefficients and near one
circle, so there are fewer than six times their least modulus.
\end{abstract}

\noindent\textbf{Keywords.} Polynomial multiples; coefficient mass; complex
roots; sectors; argument principle; Gaussian integers.

\noindent\textbf{2020 Mathematics Subject Classification.} Primary 11C08; Secondary 30C15, 26C10.

\section{Introduction}

For a nonzero $F(x)=\sum_jf_jx^j$ of degree $D$ write
\[
  \Lambda(F)=\sum_{j:f_j\ne0}\log^+|f_j|,\qquad \log^+x=\max(0,\log x),
\]
for its logarithmic coefficient mass, and $b_k(F)$ for the $k$-th largest of
the $D$ nonleading magnitudes $|f_j|$, $j<D$, counted with repetition and
zeros included ($b_k(F)=0$ for $k>D$).

For real roots
$2\le r_1\le\cdots\le r_L$, repetitions allowed, every nonzero real
multiple $F$ of $\prod_{i=1}^L(x-r_i)$ satisfies
\begin{equation}\label{eq:order}
  b_k(F)\ \ge\ |f_D|\prod_{i=k}^L(r_i-1)\quad(1\le k\le L),
  \qquad
  \Lambda(F)\ \ge\ (L+1)\log|f_D|+\sum_{i=1}^Li\log(r_i-1)
\end{equation}
\cite[Theorem~3.2 and Corollary~3.4]{coefficient-mass}, uniformly in
the degree; we call the first bound at a given $k$ the row $k$.  At complex
roots only the first row survives in general~\cite{coefficient-mass-complex}:
the polynomial $x^{2M}+\rho^{2M}$ has about $2\delta M/\pi$ roots within
$\delta$ of the imaginary axis, $b_2=0$, and mass linear in their number
\cite[Proposition~3.2]{coefficient-mass-complex}.
This paper asks what remains of \eqref{eq:order} for roots in a sector
of half-width $\delta$, where that example needs degree about $\pi K/\delta$
for $K$ roots.

\paragraph{Main results.}  For $K$ roots of modulus at least $R$ in a sector
of half-width $\delta$, multiples with leading coefficient of modulus at
least $1$ pay mass of order $K^2\log R$ when the sector is thin and the
degree small, and $(\frac\pi\delta-o(1))K\log R$, sharply, at fixed
$\delta<\pi/2$ and degree $O(K)$ as $K,R\to\infty$.  In both cases the
argument principle along a ray next to the sector produces many real zeros
of a real polynomial dominated coefficientwise by the multiple, to which
\eqref{eq:order} applies.

\paragraph{Thin sectors.}  In a sector of half-width $\delta\le\pi/4$ with
$R\ge6$ this gives mass at least $\frac{K(K-2)}8\log(\frac R2-1)$ for
multiples with leading coefficient of modulus at least $1$ and degree at
most $\pi K/(2E(2\delta))$, where
$E(2\delta)\le(8\delta/\log2)\log(e\pi/(2\delta))$ (Section~\ref{sec:thin},
Theorem~\ref{thm:thinsector}); since the degree is at least $K$, this
needs $E(2\delta)\le\pi/2$, which forces $\delta<0.033$ (and
$\delta<0.0134$ for real multiples, whose zeros off the axis come in
conjugate pairs), so the theorem is for $\delta\to0$.  At distinct
Gaussian-integer roots in the upper half-plane, where a carry argument
charges degree against mass
\cite[Theorem~6.3]{coefficient-mass-complex},
every integer multiple pays the smaller of this and
$\pi K\log(R/2)/(2E(2\delta))$ (Theorem~\ref{thm:thingauss}), which is of
order $K^2\log R$ when $\delta K\log(e\pi/(2\delta))=O(1)$; the second
term improves on the trivial bound $2K\log R$ only for
$\delta<0.0134$.

\paragraph{Sectors of every width.}  In a sector of half-width
$\delta<\pi/2$ the crossings are counted against the equidistributed share
$\delta/\pi$ of the large zeros, in the spirit of the Erd\H{o}s--Tur\'an
theorem~\cite{erdos-turan}; as $K,R\to\infty$, multiples of degree $O(K)$
with leading coefficient of modulus at least $1$ pay
$(\frac\pi\delta-o(1))K\log R$ at fixed $\delta$ and
$(\frac12-o(1))K^2\log R$ if $\delta K\to0$, both sharp
(Theorem~\ref{thm:widesector}, Corollaries~\ref{cor:wideconst}
and~\ref{cor:thinconst}, Proposition~\ref{prop:widesharp}).  At distinct
Gaussian roots in the upper half-plane and fixed $\delta$, integer
multiples of every degree pay $(\frac\pi\delta-o(1))K\log R$ as
$K,R\to\infty$ (Corollary~\ref{cor:widegauss}), and, since at most
$\pi(X+1)^2$ Gaussian integers have modulus below $X$,
$(\frac\pi\delta-o(1))K\log\max(R,\sqrt K)$ as $K\to\infty$, uniformly in
$R$ (Corollary~\ref{cor:gausscount}).  Whether $\Lambda(F)\ge cK^2\log R$
holds for all integer multiples at distinct Gaussian roots in the upper
half-plane and in a sector of fixed half-width is open.

\paragraph{Integer imitations.}  Two integer imitations of $x^N-\rho^N$ are
limited: $A((x-t)^N)$ has at most two Gaussian roots in the upper
half-plane per root of $A$ (Lemma~\ref{lem:gaussbinom}), and at odd,
pairwise coprime norms $|\alpha_j|^2$ with
$\gcd(\operatorname{Re}\alpha_j,\operatorname{Im}\alpha_j)=1$ a gap $m$ after
the lowest coefficient, that is, $m-1$ vanishing coefficients after it,
costs $2Km\log\min_j|\alpha_j|$ (Proposition~\ref{prop:gausslow}).  If
$\min_j|\alpha_j|>2$ and an integer multiple $F$ has
$b_2(F)<\min_j|\alpha_j|/2$, all the $\alpha_j$ lie in one level set
$\{y:H(y)=-v\}$, $v$ the lowest nonzero coefficient of $F$, of a polynomial
$H$ with small nonleading coefficients (Proposition~\ref{prop:gaussheavy}).
The $\alpha_j$ then lie near one circle, so $K<6\min_j|\alpha_j|$
(Proposition~\ref{prop:gausscircle}), and three Gaussian points above the
axis share such a level set at arbitrarily large $\min_j|\alpha_j|$.
Whether the number of points is bounded independently of
$\min_j|\alpha_j|$ is open.

\paragraph{Conventions.}
All logarithms are natural.  For $F\in\mathbb C[x]$ of degree $D$ we write
$F^*(t)=t^DF(1/t)=\sum_{d=0}^Df_{D-d}t^d$ for the reciprocal polynomial.
Coefficients of a polynomial named by a capital letter, when used, carry the
lowercase letter: $F=\sum_jf_jx^j$, $B=\sum_ib_ix^i$.  With an argument,
$b_k(\cdot)$ is always the $k$-th largest nonleading magnitude, so $b_k(B)$
and the coefficient $b_k$ of $B$ differ.
References of the form \cite[Theorem~3.2]{coefficient-mass} are to the
first companion paper, and those of the form
\cite[Theorem~6.3]{coefficient-mass-complex}
to the second.

\section{Thin sectors}\label{sec:thin}

This section proves that $K$ roots of modulus at least $R\ge6$ in a sector
of half-width $\delta\le\pi/4$ force mass at least
$\frac{K(K-2)}8\log(\frac R2-1)$ on every multiple with leading coefficient
of modulus at least $1$ and degree at most $\pi K/(2E(2\delta))$, with $E$
defined below (Theorem~\ref{thm:thinsector}).  Without separated moduli, a multiple may have a single tropical root
\cite[Theorem~3.3]{coefficient-mass-complex} for all the prescribed roots of an annulus, as $x^{2M}+\rho^{2M}$ and
$x^{2m}+x+a$ of \cite[Proposition~3.2]{coefficient-mass-complex} do.  The angle mechanism of
\cite[Proposition~3.2]{coefficient-mass-complex} puts $K$ roots within an angle $\delta$ of a
ray only at the price of degree about $\pi K/\delta$.  A multiple of smaller
degree is handled by the real theory, applied to the real polynomial
$r\mapsto\operatorname{Im}\bigl(e^{-i\theta}F(re^{i\psi})\bigr)$ along a ray
$\arg z=\psi$ next to the sector.  For $0<x\le\pi/2$ put
\[
  E(x)=\frac{4x}{\log2}\Bigl(1+\log\Bigl(1+\frac{\pi}{2x}\Bigr)\Bigr).
\]
$E$ is increasing, since
$\frac{d}{dx}\bigl(x(1+\log(1+\frac{\pi}{2x}))\bigr)=1+\log(1+t)-\frac t{1+t}>0$
with $t=\pi/(2x)$.

\begin{lemma}[Crossings next to a sector]\label{lem:sectorcount}
Let $B\in\mathbb C[x]$ be nonzero of degree $d$, let $\varphi\in\R$,
$0<\delta'\le\pi/2$ and $R>0$, and suppose that $B$ has no zero other than $0$
on the two rays $\arg z=\varphi\pm\delta'$, and at least $K$ zeros, counted
with multiplicity, in $\{z:|z|\ge R,\ |\arg z-\varphi|<\delta'\}$.  Put
$E(x)=\frac{4x}{\log2}\bigl(1+\log\bigl(1+\frac{\pi}{2x}\bigr)\bigr)$ for
$0<x\le\pi/2$.  Then there are
$r_1\in[R/2,R)$ and $\psi\in\{\varphi-\delta',\varphi+\delta'\}$ such that for
every $\theta\in\R$ the real polynomial
\[
  h_\theta(r)=\operatorname{Im}\bigl(e^{-i\theta}B(re^{i\psi})\bigr)
  =\sum_{k=0}^d\operatorname{Im}\bigl(e^{-i\theta}b_ke^{ik\psi}\bigr)\,r^k
\]
has at least $K-dE(\delta')/\pi-1$ distinct zeros in $(r_1,\infty)$.
\end{lemma}

\begin{proof}
For $r>0$ let $\Gamma_r=\{re^{it}:|t-\varphi|\le\delta'\}$, traversed in the
direction of increasing $t$, and for a point $\beta$ with $|\beta|\ne r$ let
$\Delta_\beta(r)$ be the change of a continuous argument of $z-\beta$ along
$\Gamma_r$.  We claim
\[
  |\Delta_\beta(r)|\ \le\ \min\Bigl(2\pi,\ \frac{2\delta'r}{\bigl|r-|\beta|\bigr|}\Bigr).
\]
The second bound holds because $|d\arg(z-\beta)|\le|dz|/|z-\beta|$, the arc
has length $2\delta'r$, and $|z-\beta|\ge|r-|\beta||$ on it.  For the first:
if $|\beta|<r$, the argument of $re^{it}-\beta$ has derivative
$(r^2-\operatorname{Re}(\bar\beta re^{it}))/|re^{it}-\beta|^2>0$ and increases by
$2\pi$ over a full turn, so $0<\Delta_\beta(r)<2\pi$; if $|\beta|>r$, the
points $z-\beta$, $|z|=r$, lie in a closed disc not containing $0$, which is
seen from $0$ under an angle less than $\pi$, so $|\Delta_\beta(r)|<\pi$.

Next, for every $\varrho\ge0$ and $a>0$,
\[
  \int_a^{2a}\min\Bigl(2\pi,\frac{2\delta'r}{|r-\varrho|}\Bigr)\frac{dr}r
  \ \le\ 4\delta'\Bigl(1+\log\Bigl(1+\frac\pi{2\delta'}\Bigr)\Bigr)
  =E(\delta')\log2 .
\]
Indeed, the integrand is $\min(2\pi/r,2\delta'/|r-\varrho|)$.  Put
$\eta=\delta'a/\pi$.  Where $|r-\varrho|\le\eta$ it is at most $2\pi/a$ on a
set of measure at most $2\eta$, which contributes at most $4\delta'$.  Elsewhere
it is at most $2\delta'/|r-\varrho|$ on a set of measure at most $a$ on which
$|r-\varrho|>\eta$; such an integral is largest for the set
$\eta<|r-\varrho|\le\eta+a/2$, where it equals
$4\delta'\log(1+a/(2\eta))=4\delta'\log(1+\pi/(2\delta'))$.

Let $\beta_1,\ldots,\beta_d$ be the zeros of $B$ with multiplicity, and
$S(r)=\sum_l\min(2\pi,2\delta'r/|r-|\beta_l||)$.  By the last display,
$\int_a^{2a}S(r)\,dr/r\le dE(\delta')\log2$, while
$\int_a^{2a}dr/r=\log2$; so $S(r)\le dE(\delta')$ on a subset of $[a,2a]$ of
positive measure.  Taking $a=R/2$, choose $r_1\in[R/2,R)$ there, different from
every $|\beta_l|$; taking $a=M$ larger than every $|\beta_l|$, choose $r_2$ in
the same way in $[M,2M]$.  Since
$\arg B(z)=\arg b_d+\sum_l\arg(z-\beta_l)$, the change $\Delta_{\rm arc}(r_i)$
of $\arg B$ along $\Gamma_{r_i}$ satisfies
$|\Delta_{\rm arc}(r_i)|\le S(r_i)\le dE(\delta')$.

Let $\psi_\pm=\varphi\pm\delta'$ and let $\Delta_\pm$ be the change of
$\arg B(re^{i\psi_\pm})$ as $r$ increases from $r_1$ to $r_2$.  The region
$W=\{r_1<|z|<r_2,\ |\arg z-\varphi|<\delta'\}$ has no zero of $B$ on its
boundary and contains the $K$ zeros of the statement (their moduli are at
least $R>r_1$ and below $M<r_2$).  The argument principle along
$\partial W$, traversed counterclockwise, gives
\[
  2\pi K\ \le\ \Delta_-+\Delta_{\rm arc}(r_2)-\Delta_+-\Delta_{\rm arc}(r_1)
  \ \le\ |\Delta_-|+|\Delta_+|+2dE(\delta').
\]
Let $\psi$ be the ray with the larger $|\Delta_\pm|$, so that
$|\Delta_\psi|\ge\pi K-dE(\delta')$, and let $A(r)$ be a continuous argument of
$B(re^{i\psi})$ on $[r_1,r_2]$.  Fix $\theta$.  The open interval between
$A(r_1)$ and $A(r_2)$ has length $|\Delta_\psi|$, so it contains at least
$|\Delta_\psi|/\pi-1$ points of $\theta+\pi\mathbb Z$; by continuity each of
them equals $A(r)$ for some $r\in(r_1,r_2)$, and distinct points give distinct
$r$.  At such $r$ the number $e^{-i\theta}B(re^{i\psi})$ is real, so
$h_\theta(r)=0$.  This gives at least $|\Delta_\psi|/\pi-1\ge K-dE(\delta')/\pi-1$
distinct zeros of $h_\theta$ in $(r_1,r_2)$.
\end{proof}

\begin{theorem}[Thin sectors]\label{thm:thinsector}
Let $B\in\mathbb C[x]$ have degree $d$ and $|b_d|\ge1$, let $\varphi\in\R$,
$0<\delta\le\pi/4$ and $R\ge6$, and suppose that $B$ has at least $K$ zeros,
counted with multiplicity, in $\{z:|z|\ge R,\ |\arg z-\varphi|\le\delta\}$.  If
$d\,E(2\delta)\le\pi K/2$, with $E$ as in Lemma~\ref{lem:sectorcount}, then
with $m=\lceil K/2\rceil-1$,
\[
  b_k(B)\ \ge\ |b_d|\Bigl(\frac R2-1\Bigr)^{m-k+1}\quad(1\le k\le m),
\]
\[
  \Lambda(B)\ \ge\ (m+1)\log|b_d|+\frac{m(m+1)}2\log\Bigl(\frac R2-1\Bigr),
\]
and $\frac{m(m+1)}2\ge\frac{K(K-2)}8$.
\end{theorem}

Since $d\ge K$, the degree hypothesis can hold only when
$E(2\delta)\le\pi/2$, which forces $\delta<0.033$ (at $\delta=0.033$,
$8\delta/\log2>0.38$ and $1+\log(1+\pi/(4\delta))>1+\log24.7>4.2$, so
$E(2\delta)>1.596>\pi/2$, and $E$ is increasing; numerically the threshold
is $\delta=0.0321\ldots$).  If $B$ is real and the sector misses the real
axis, the conjugates of the $K$ zeros give $d\ge2K$, so the hypothesis needs
$E(2\delta)\le\pi/4$, that is $\delta<0.0134$: the theorem is meant for
$\delta\to0$, and Section~\ref{sec:widesectors} treats wider sectors.

\begin{proof}
Choose $\delta'\in(\delta,2\delta]$ such that no zero of $B$ other than $0$ lies on the rays
$\arg z=\varphi\pm\delta'$; all but finitely many values qualify.  The $K$
zeros lie in the open sector $|\arg z-\varphi|<\delta'$, and $E(\delta')\le
E(2\delta)$ because $E$ is increasing.  Lemma~\ref{lem:sectorcount} gives
$r_1\ge R/2$ and a ray $\psi$ such that every $h_\theta$ has at least
$K-dE(2\delta)/\pi-1\ge K/2-1$ distinct zeros in $(r_1,\infty)$, hence at least
$m$ of them, say $\varrho_1<\cdots<\varrho_m$, all larger than $R/2\ge3$.
Choose $\theta$ with $e^{-i\theta}b_de^{id\psi}=i|b_d|$.  Then $h_\theta$ is a
nonzero real polynomial of degree $d$ with leading coefficient $|b_d|$ and
coefficients $h_k=\operatorname{Im}(e^{-i\theta}b_ke^{ik\psi})$, so
$|h_k|\le|b_k|$, and it is divisible by $\prod_{i\le m}(x-\varrho_i)$ with
$\varrho_1\ge2$.  \cite[Theorem~3.2]{coefficient-mass} and \cite[Corollary~3.4]{coefficient-mass} give
$b_k(h_\theta)\ge|b_d|\prod_{i=k}^m(\varrho_i-1)$ and
$\Lambda(h_\theta)\ge(m+1)\log|b_d|+\sum_{i=1}^mi\log(\varrho_i-1)$, and
$\varrho_i-1>R/2-1$.  The coefficientwise domination $|h_k|\le|b_k|$, with the
leading coefficients at the same position $d$, gives $b_k(B)\ge b_k(h_\theta)$
and $\Lambda(B)\ge\Lambda(h_\theta)$.  Finally $m+1\ge K/2$ and $m\ge K/2-1$.
\end{proof}

The theorem uses no integrality, and its degree hypothesis cannot be
dropped: $x^{2M}+\rho^{2M}$ in \cite[Proposition~3.2]{coefficient-mass-complex} has about
$2\delta M/\pi$ roots within $\delta$ of the positive imaginary axis and
$b_2=0$, and its degree $2M$ exceeds $\pi K/(2E(2\delta))$.  At Gaussian roots
the large-degree case is paid for by integrality, as the next section shows.

\section{Gaussian-integer roots}\label{sec:gaussian}

For roots in $\mathbb Z[i]$ and integer multiples there is one more piece of
information: an integer polynomial takes Gaussian-integer values at them,
and a nonzero Gaussian integer has modulus at least one.  This charges the
degree of a multiple against its mass~\cite{coefficient-mass-complex}, which
replaces the degree hypothesis of Theorem~\ref{thm:thinsector} by the
alternative bound $\pi K\log(R/2)/(2E(2\delta))$ (Theorem~\ref{thm:thingauss}).

Throughout, $\alpha_1,\ldots,\alpha_K\in\mathbb Z[i]$ are distinct, with
$\alpha_j=a_j+ic_j$ and $c_j\ge1$; put $\rho_j=|\alpha_j|$,
$\rho_{\min}=\min_j\rho_j$ and
\[
  P=\prod_{j=1}^K(x-\alpha_j)(x-\bar\alpha_j)
   =\prod_{j=1}^K\bigl(x^2-2a_jx+a_j^2+c_j^2\bigr)\in\mathbb Z[x],
\]
a monic polynomial with $2K$ distinct roots and $P(0)\ne0$.  For
$F=\sum_{i=0}^Df_ix^i$ and $0\le m\le D+1$ write
\[
  F_{<m}=\sum_{i<m}f_ix^i,\qquad G_m=\sum_{i\ge m}f_ix^{i-m},\qquad
  F=F_{<m}+x^mG_m .
\]
We call $m\in\{1,\ldots,D\}$ a \emph{split} of $F$.  A split is \emph{clean
at} $\alpha$ if $G_m(\alpha)=0$, and \emph{$P$-clean} if it is clean at every
$\alpha_j$.  A polynomial $B\in\mathbb Z[x]$ is \emph{$P$-irreducible} if
$B(0)\ne0$ and no split of $B$ is $P$-clean.

We use three results of~\cite{coefficient-mass-complex} in this setting.
Let $F\in\mathbb Z[x]$ have degree $D$, let $\alpha\in\mathbb Z[i]$ with
$F(\alpha)=0$ and $\rho=|\alpha|>2$, and let $m$ be a split of $F$ that is
not clean at $\alpha$.  Then some $i<m$ has $|f_i|\ge(\rho/2)^{m-i}$
\cite[Lemma~6.1]{coefficient-mass-complex}.
Every nonzero $F\in\mathbb Z[x]$ divisible by $P$ can be written
\[
  F=\sum_{s=1}^rx^{e_s}B_s,\qquad e_s+\deg B_s<e_{s+1},
\]
with every $B_s\in\mathbb Z[x]$ divisible by $P$ and $P$-irreducible; then
$\Lambda(F)=\sum_s\Lambda(B_s)$, $|\operatorname{lead}B_s|\ge1$, and
$b_k(F)\ge b_k(B_s)$ for every $s$ and $k$
\cite[Lemma~6.2]{coefficient-mass-complex}.
We call the $B_s$ the blocks of the decomposition.  Finally, if
$\rho_{\min}>2$, every $P$-irreducible $B\in\mathbb Z[x]$ divisible by $P$,
of degree $d$, has
\[
  \sum_{i<d}\log^+|b_i|\ \ge\ d\log\frac{\rho_{\min}}2,
  \qquad\text{so}\qquad
  \Lambda(B)\ \ge\ \log|b_d|+d\log\frac{\rho_{\min}}2
\]
\cite[Theorem~6.3]{coefficient-mass-complex}.

\begin{theorem}[Thin Gaussian clusters]\label{thm:thingauss}
In the setting of Section~\ref{sec:gaussian}, let $\varphi\in\R$,
$0<\delta\le\pi/4$ and $R\ge6$, and suppose that every
$\alpha_j$ has $|\alpha_j|\ge R$ and $|\arg\alpha_j-\varphi|\le\delta$.  With
$E$ as in Lemma~\ref{lem:sectorcount}, every
nonzero $F\in\mathbb Z[x]$ divisible by $P$ satisfies $\Lambda(F)\ge2K\log R$
and
\[
  \Lambda(F)\ \ge\ \min\Bigl\{\frac{K(K-2)}8\log\Bigl(\frac R2-1\Bigr),\
  \frac{\pi K\log(R/2)}{2E(2\delta)}\Bigr\},
  \qquad
  \frac{\pi}{2E(2\delta)}\ \ge\ \frac{\pi\log2}{16\,\delta\log(e\pi/(2\delta))}.
\]
If moreover some block $B$ of some decomposition of $F$ as in
\cite[Lemma~6.2]{coefficient-mass-complex} has
$\deg B\cdot E(2\delta)\le\pi K/2$, then
$b_k(F)\ge(R/2-1)^{m-k+1}$ for $1\le k\le m=\lceil K/2\rceil-1$.
\end{theorem}

\begin{proof}
Let $B$ be a block, of degree $d$, of any decomposition of $F$ as in
\cite[Lemma~6.2]{coefficient-mass-complex}; for the rows, take the block of the hypothesis.
It is divisible by $P$, so it
vanishes at the $K$ distinct $\alpha_j$; $B(0)\ne0$, $|b_d|\ge1$,
$\Lambda(F)\ge\Lambda(B)$ and $b_k(F)\ge b_k(B)$.  If $dE(2\delta)\le\pi K/2$,
Theorem~\ref{thm:thinsector} applies to $B$ and gives the first term, and the
rows.  Otherwise $d>\pi K/(2E(2\delta))$, and \cite[Theorem~6.3]{coefficient-mass-complex},
with $\rho_{\min}\ge R>2$, gives
$\Lambda(B)\ge d\log(\rho_{\min}/2)>\pi K\log(R/2)/(2E(2\delta))$.  Next,
$B=PQ$ with $Q\in\mathbb Z[x]$ since $P$ is monic, so $B(0)=P(0)Q(0)$ is a
nonzero multiple of $P(0)=\prod_j\rho_j^2\ge R^{2K}$, and $B(0)$ is a
coefficient of $F$; this gives $\Lambda(F)\ge2K\log R$.  Finally, for
$0<\delta\le\pi/4$ we have $1+\pi/(4\delta)\le\pi/(2\delta)$, so
$E(2\delta)\le(8\delta/\log2)(1+\log(\pi/(2\delta)))
=(8\delta/\log2)\log(e\pi/(2\delta))$.
\end{proof}

\section{Sectors of every width}\label{sec:widesectors}

This section proves that in a sector of fixed half-width $\delta<\pi/2$,
every multiple of degree $O(K)$ with $|f_D|\ge1$ at $K$ roots of modulus at
least $R$ pays $(\pi/\delta-o(1))K\log R$ as $K,R\to\infty$, that this
constant is attained, and that at Gaussian roots the degree restriction can
be removed.  The loss $dE(\delta')/\pi$ in Lemma~\ref{lem:sectorcount} is what confines
Theorem~\ref{thm:thinsector} to thin sectors; here it is replaced by
$\delta'n/\pi$, where $n$ counts only the zeros of modulus above the inner
arc, plus a degree term that becomes small when the inner radius is averaged
over a long range of scales.  The rate $\delta'/\pi$ is the share of a
sector of opening $2\delta'$ in a full turn, and $x^N-\rho^N$ shows that no
smaller rate is possible.  $K$ zeros of modulus at least $R$ in a sector of
half-width $\delta<\pi/2$ then cost either about $\frac12m^2\log R$, where
$m$ is the excess of $K$ over the share $\delta n/\pi$, or about $n\log R$
through the Mahler measure, which gives the constant $\pi/\delta$
(Theorem~\ref{thm:widesector}, Corollary~\ref{cor:wideconst},
Proposition~\ref{prop:widesharp}).  At Gaussian roots and fixed $\delta$ the
degree restriction is removed by \cite[Theorem~6.3]{coefficient-mass-complex}
(Corollary~\ref{cor:widegauss}), and every integer multiple pays
$(\frac\pi{2\delta}-o(1))K\log K$ as $K\to\infty$, whatever $R$
(Corollary~\ref{cor:gausscount}).
In wide sectors a bound $\Lambda\ge cK^2\log R$ valid for all root sets is
therefore false for real multiples.  For integer
multiples at Gaussian roots it is open; Section~\ref{sec:imitations} gives two
limits on integer imitations of $x^N-\rho^N$ (Lemma~\ref{lem:gaussbinom},
Proposition~\ref{prop:gausslow}) and a reduction of the second row to level
sets (Propositions~\ref{prop:gaussheavy} and~\ref{prop:gausscircle}).

For $0<\delta\le\pi/2$ put
\[
  \kappa(\delta)=\frac1\pi\int_0^1\min\Bigl(2\arcsin q,\ \frac{2\delta q}{1-q}\Bigr)\frac{dq}q .
\]
The integrand increases with $\delta$ and is at most $2\arcsin(q)/q$, so
$\kappa$ is increasing and, by dominated convergence, continuous.  Since
$\int_0^1\arcsin(q)\,dq/q=\int_0^{\pi/2}s\cot s\,ds=-\int_0^{\pi/2}\log\sin s\,ds
=\frac\pi2\log2$, we have $\kappa(\delta)\le\log2$; and bounding
$2\arcsin q\le\pi q$ and splitting at $q=1-2\delta/\pi$,
\[
  \kappa(\delta)\ \le\ \frac1\pi\Bigl(\int_0^{1-2\delta/\pi}\frac{2\delta\,dq}{1-q}
  +\int_{1-2\delta/\pi}^1\pi\,dq\Bigr)
  =\frac{2\delta}\pi\Bigl(1+\log\frac\pi{2\delta}\Bigr).
\]

\begin{lemma}[Sharp crossings next to a sector]\label{lem:sharpcross}
Let $B\in\mathbb C[x]$ be nonzero of degree $d$, let $\varphi\in\R$ and
$0<\delta'\le\pi/2$, and suppose that $B$ has no zero other than $0$ on the
rays $\arg z=\varphi\pm\delta'$.  Let $0<a<b$ and
$\varepsilon=\kappa(\delta')d/\log(b/a)$.  There is $r_1\in[a,b)$, different
from the modulus of every zero of $B$, with the following property.  Let $n$
be the number of zeros of $B$ of modulus greater than $r_1$, and $N$ the
number of those in the open sector $|\arg z-\varphi|<\delta'$, both counted
with multiplicity.  For $\psi\in\R$ and $\theta\in\R$ put
$h_\theta(r)=\operatorname{Im}\bigl(e^{-i\theta}B(re^{i\psi})\bigr)$.
There is $\psi\in\{\varphi-\delta',\varphi+\delta'\}$ such that for
every $\theta$ the polynomial $h_\theta$ has at least
$N-\delta'n/\pi-\varepsilon-1$ distinct zeros in $(r_1,\infty)$.
\end{lemma}

\begin{proof}
For $r>0$ let $\Gamma_r$ and $\Delta_\beta(r)$ be as in the proof of
Lemma~\ref{lem:sectorcount}, let $[\,\cdot\,]$ be the indicator of a
condition, and for $\beta\ne0$ with $|\beta|\ne r$ put
$q=\min(|\beta|/r,r/|\beta|)<1$.  We claim
\[
  \bigl|\Delta_\beta(r)-2\delta'[\,|\beta|<r\,]\bigr|\ \le\
  c(q)=\min\Bigl(2\arcsin q,\ \frac{2\delta'q}{1-q}\Bigr).
\]
If $|\beta|<r$, write $re^{it}-\beta=re^{it}(1-w)$ with
$w=(\beta/r)e^{-it}$; if $|\beta|>r$, write $re^{it}-\beta=-\beta(1-w)$ with
$w=(r/\beta)e^{it}$.  In both cases $|w|=q<1$, so $1-w$ stays in the disc of
radius $q$ about $1$, where the principal argument is continuous and at most
$\arcsin q$ in absolute value.  So $\Delta_\beta(r)-2\delta'[\,|\beta|<r\,]$
is the change of $\operatorname{Arg}(1-w)$ along $\Gamma_r$, which is at most
$2\arcsin q$ in absolute value; and since $dw/dt=\mp iw$, the derivative
$\frac{d}{dt}\operatorname{Arg}(1-w)=\operatorname{Im}\frac{\pm iw}{1-w}$ is at most $q/(1-q)$ in
absolute value, over a range of length $2\delta'$.  For $\beta=0$,
$\Delta_0(r)=2\delta'$ exactly.

Let $\beta_1,\ldots,\beta_d$ be the zeros of $B$ with multiplicity (all but
the zero at $0$ are in the sum below) and
$S(r)=\sum_{\beta_l\ne0}c(q_l(r))$.  For each $\beta_l\ne0$,
\[
  \int_{-\infty}^{\infty}c\bigl(q_l(e^u)\bigr)\,du
  =2\int_0^\infty c(e^{-v})\,dv=2\int_0^1c(q)\frac{dq}q=2\pi\kappa(\delta'),
\]
so $\int_{\log a}^{\log b}S(e^u)\,du\le2\pi\kappa(\delta')d$ and
$S(e^u)\le2\pi\varepsilon$ on a subset of $[\log a,\log b)$ of positive
measure.  Choose $r_1=e^u$ there, different from every $|\beta_l|$.

Put $X_1=N-\delta'n/\pi-\varepsilon-1$ and choose $\eta>0$
with $\lceil X_1-\eta\rceil=\lceil X_1\rceil$.  Since $c(q)\le\pi q$, $S(r)\to0$ as
$r\to\infty$; choose $r_2$ larger than $r_1$ and every $|\beta_l|$ with
$S(r_2)\le2\pi\eta$.  Let $\psi_\pm=\varphi\pm\delta'$ and let $\Delta_\pm$ be
the change of a continuous argument of $B(re^{i\psi_\pm})$ as $r$ increases
from $r_1$ to $r_2$.  The region
$W=\{r_1<|z|<r_2,\ |\arg z-\varphi|<\delta'\}$ has no zero of $B$ on its
boundary and contains exactly $N$ zeros.  With $\Delta_{\rm arc}(r)$ as in
the proof of Lemma~\ref{lem:sectorcount}, the argument principle gives
$2\pi N=\Delta_-+\Delta_{\rm arc}(r_2)-\Delta_+-\Delta_{\rm arc}(r_1)$.  By
the claim,
$\Delta_{\rm arc}(r)=\sum_l\Delta_{\beta_l}(r)=2\delta'\#\{l:|\beta_l|<r\}+\gamma(r)$
with $|\gamma(r)|\le S(r)$, and
$\#\{l:|\beta_l|<r_2\}-\#\{l:|\beta_l|<r_1\}=n$.  Hence
\[
  |\Delta_-|+|\Delta_+|\ \ge\ |\Delta_--\Delta_+|
  \ \ge\ 2\pi N-2\delta'n-S(r_1)-S(r_2)
  \ \ge\ 2\pi N-2\delta'n-2\pi\varepsilon-2\pi\eta .
\]
By the last paragraph of the proof of Lemma~\ref{lem:sectorcount}, for every
$\theta'$ there are at least $|\Delta_\pm|/\pi-1$ distinct
$r\in(r_1,r_2)$ at which $e^{-i\theta'}B(re^{i\psi_\pm})$ is real.

Let $\psi$ be the ray with the larger $|\Delta_\pm|$.  Then
$|\Delta_\psi|/\pi-1\ge X_1-\eta$, so every $h_\theta$ has at least
$\lceil X_1-\eta\rceil=\lceil X_1\rceil\ge X_1$ distinct zeros in
$(r_1,\infty)$.
\end{proof}

\begin{theorem}[Sectors of every width]\label{thm:widesector}
Let $B\in\mathbb C[x]$ have degree $d$ and $|b_d|\ge1$, let $\varphi\in\R$,
$0<\delta<\pi/2$, $T>1$ and $R\ge3T$, and suppose that $B$ has at least $K$
zeros, counted with multiplicity, in
$\{z:|z|\ge R,\ |\arg z-\varphi|\le\delta\}$.  Put
\[
  \kappa(\delta)=\frac1\pi\int_0^1\min\Bigl(2\arcsin q,\ \frac{2\delta q}{1-q}\Bigr)\frac{dq}q,
\]
so that $\kappa(\delta)\le\log2$ and
$\kappa(\delta)\le\frac{2\delta}\pi(1+\log\frac\pi{2\delta})$,
$\omega=\kappa(\delta)d/\log T$ and, for $\nu\ge0$,
$g(\nu)=\nu\log(R/T)-\log(\nu+1)-1$.  For every integer $m\ge0$,
\[
  \Lambda(B)\ \ge\ \min\Bigl\{\frac{m(m+1)}2\log\Bigl(\frac RT-1\Bigr),\
  g(\nu_m)\Bigr\},\qquad
  \nu_m=\frac\pi\delta\,(K-1-m-\omega)^+ .
\]
\end{theorem}

\begin{proof}
Choose $\delta'\in(\delta,\pi/2)$ such that no zero of $B$ other than $0$ lies
on the rays $\arg z=\varphi\pm\delta'$; all but finitely many values qualify.
Apply Lemma~\ref{lem:sharpcross} with $a=R/T$ and $b=R$, so
$\varepsilon'=\kappa(\delta')d/\log T$; the $K$ zeros have modulus at least
$R>r_1$ and lie in the open sector, so $N\ge K$.  Note $r_1\ge R/T\ge3$, and
$g$ is increasing on $[0,\infty)$, as $g'(\nu)=\log(R/T)-1/(\nu+1)>0$.

\emph{Mahler measure.}  Suppose $n\ge1$ and let
$B^*(t)=t^dB(1/t)=\sum_kb_{d-k}t^k$, so $B^*(0)=b_d$ and the nonzero zeros of
$B^*$ are the $1/\beta$, $\beta\ne0$ a zero of $B$.  Jensen's formula on
$|t|=\tau=n/(n+1)\ge\frac12$, where the $n$ points $1/\beta$ with
$|\beta|>r_1$ lie in $|t|<1/3$ and each contributes $\log(\tau|\beta|)\ge\log(\tau r_1)$,
and every other contribution is nonnegative, gives
\[
  \log\max_{|t|=\tau}|B^*(t)|\ \ge\ \log|b_d|+n\log(\tau r_1)\ \ge\ n\log r_1-1,
\]
using $n\log\frac{n+1}n\le1$.  On the other side
$\max_{|t|=\tau}|B^*|\le\max_i|b_i|\sum_{k\ge0}\tau^k=(n+1)\max_i|b_i|$.  So
$\Lambda(B)\ge\log\max_i|b_i|\ge n\log r_1-\log(n+1)-1\ge g(n)$.

\emph{Crossings.}  Choose $\theta$ with $e^{-i\theta}b_de^{id\psi}=i|b_d|$.  As
in the proof of Theorem~\ref{thm:thinsector}, $h=h_\theta$ is real of degree
$d$ with leading coefficient $|b_d|\ge1$ and $|h_k|\le|b_k|$, so
$\Lambda(B)\ge\Lambda(h)$.

Let $X_1=N-\delta'n/\pi-\varepsilon'-1$.  If $X_1\ge m$, then by
Lemma~\ref{lem:sharpcross} $h$ has $m$ distinct zeros
$\varrho_1<\cdots<\varrho_m$ larger than $r_1\ge3$, and \cite[Theorem~3.2]{coefficient-mass} with \cite[Corollary~3.4]{coefficient-mass} gives
$\Lambda(h)\ge\sum_{i=1}^mi\log(\varrho_i-1)\ge\frac{m(m+1)}2\log(R/T-1)$.
Otherwise $n>(\pi/\delta')(K-1-m-\varepsilon')$.  If the right side is positive,
$n\ge1$ and $\Lambda(B)\ge g(n)\ge g(\nu_m(\delta'))$, where $\nu_m(\delta')$
is $\nu_m$ with $\delta'$ and $\varepsilon'$ in place of $\delta$ and $\omega$; if not,
$\nu_m(\delta')=0$ and $g(0)=-1<0\le\Lambda(B)$.  So the bound holds with
$\nu_m(\delta')$, for all admissible $\delta'$ arbitrarily close to $\delta$;
as $\delta'\downarrow\delta$, $\kappa(\delta')\to\kappa(\delta)$ and
$\nu_m(\delta')\to\nu_m$, and $g$ is continuous.

\end{proof}

The first term is the thin-sector mechanism, now at every width; the second
is the angle mechanism.  If the zeros of $B$ of modulus above $R/T$ are few,
for instance if $B=PQ$ with $P$ the product of $K$ conjugate pairs
$\beta_j,\bar\beta_j$ with $\beta_j$ in the sector and all zeros of $Q$ in $|z|\le R/T$, then $n\le2K$ in the proof.
For $m$ with $\nu_m>2K$ its alternative $n>(\pi/\delta')(K-1-m-\varepsilon')$
fails for $\delta'$ near $\delta$, so the crossings step applies; the largest
such $m$ is at least $(1-\frac{2\delta}\pi)K-\omega-2$, and whenever this is
nonnegative,
$\Lambda(B)\ge\frac12\bigl((1-\frac{2\delta}\pi)K-\omega-2\bigr)^2\log(R/T-1)$.
So in a sector of any
width short of a half-plane a cheap multiple needs about
$(\frac\pi\delta-2)K$ further zeros of modulus above $R/T$, which the
Mahler term charges linearly.

\begin{corollary}[The constant $\pi/\delta$]\label{cor:wideconst}
In Theorem~\ref{thm:widesector} suppose $d\le\lambda K$ with $\lambda>0$, let
$0<\varepsilon\le\frac12$, and take $T=\exp(\kappa(\delta)\lambda/\varepsilon)$,
assuming $R\ge3T$.  Then
\[
  \Lambda(B)\ \ge\ \min\Bigl\{\frac{\varepsilon^2K^2}2\log\Bigl(\frac RT-1\Bigr),\
  \frac\pi\delta\bigl((1-2\varepsilon)K-2\bigr)\log\frac RT
  -\log\Bigl(\frac{\pi K}\delta+1\Bigr)-1\Bigr\}.
\]
Consequently, for fixed $\delta$ and $\lambda$,
$\Lambda(B)\ge(\pi/\delta-o(1))K\log R$ as $K\to\infty$ and $R\to\infty$.
\end{corollary}

\begin{proof}
Here $\omega\le\kappa(\delta)\lambda K/\log T=\varepsilon K$.  Take
$m=\lceil\varepsilon K\rceil\le\varepsilon K+1$: then
$\frac{m(m+1)}2\ge\frac{\varepsilon^2K^2}2$ and
$K-1-m-\omega\ge(1-2\varepsilon)K-2=:\nu_0\delta/\pi$.  Since $g$ is increasing,
$g(\nu_m)\ge g(\nu_0^+)$, and $g(\nu_0)\ge\nu_0\log(R/T)-\log(\pi K/\delta+1)-1$
because $\nu_0\le\pi K/\delta$; if $\nu_0<0$ the claimed bound is negative.
For the limit take $\varepsilon=\max(K^{-1/3},(\log R)^{-1/2})$: then
$\log T=O((\log R)^{1/2})$, the first term is at least
$\frac12K^{4/3}\log(R/T-1)$, which exceeds the second for large $K$, and the second is
$(\pi/\delta)(1-o(1))K\log R$.
\end{proof}

\begin{proposition}[Sharpness of $\pi/\delta$]\label{prop:widesharp}
Let $\varphi\in\R$, $0<\delta\le\pi$, $K\ge1$, $\rho\ge1$ and
$N=\lceil\pi K/\delta\rceil$.  The monic polynomial $x^N-\rho^N$ has at least
$K$ zeros in $\{z:|z|=\rho,\ |\arg z-\varphi|\le\delta\}$, degree
$N\le\pi K/\delta+1$, and
\[
  \Lambda(x^N-\rho^N)=N\log\rho\ \le\ \Bigl(\frac{\pi K}\delta+1\Bigr)\log\rho .
\]
It is real, and integral when $\rho^N\in\mathbb Z$.  If $0<\varphi-\delta$
and $\varphi+\delta<\pi$, it is a real multiple of $K$ conjugate pairs
$\beta_j,\bar\beta_j$ with $\beta_j$ in the sector.
\end{proposition}

\begin{proof}
The zeros $\rho e^{2\pi il/N}$ have arguments spaced by $2\pi/N$, so a closed
arc of length $2\delta$ contains at least $\lfloor\delta N/\pi\rfloor\ge K$ of
them.  The only nonleading nonzero coefficient is $-\rho^N$.  Under the last
hypothesis the $K$ zeros are nonreal, and their conjugates are zeros of the
real polynomial $x^N-\rho^N$.
\end{proof}

Its degree is at most $(\pi/\delta+1)K$, so Corollary~\ref{cor:wideconst}
applies with $\lambda=\pi/\delta+1$: the constant $\pi/\delta$ of the angle
mechanism is sharp in every sector of half-width $\delta<\pi/2$.  In a sector of fixed half-width no
lower bound valid for all sets of $K$ roots is therefore better than linear
in $K$ (particular root sets, such as roots on one ray, still cost
quadratic mass by \cite[Corollary~2.4]{coefficient-mass-complex} with
\cite[Corollary~3.4]{coefficient-mass}, applied to $F(e^{i\varphi}x)$ when
$|f_D|\ge1$), and $\Lambda\ge cK^2\log R$ fails already
for the monic real multiple $x^N-\rho^N$.  The thin-sector term is sharp too.

\begin{corollary}[The constant $\frac12$]\label{cor:thinconst}
In Theorem~\ref{thm:widesector} suppose $d\le\lambda K$ with $\lambda>0$, and
let $m=K-1-\lceil\omega+\delta K^2/(2\pi)\rceil$.  If $m\ge0$, then
\[
  \Lambda(B)\ \ge\ \min\Bigl\{\frac{m(m+1)}2\log\Bigl(\frac RT-1\Bigr),\
  \frac{K^2}2\log\frac RT-\log\Bigl(\frac{K^2}2+1\Bigr)-1\Bigr\}.
\]
Consequently, for fixed $\lambda$ and $T$, $\Lambda(B)\ge(\frac12-o(1))K^2\log R$
as $K\to\infty$, $R\to\infty$ and $\delta K\to0$.  Conversely, for
$\rho_1,\ldots,\rho_K\in[R,2R]$ the monic multiple
$\prod_j(x-\rho_je^{i\varphi})$ of degree $K$ has
$\Lambda\le\frac{K(K+1)}2\log(2R)+K^2\log2$.
\end{corollary}

\begin{proof}
$K-1-m-\omega=\lceil\omega+\delta K^2/(2\pi)\rceil-\omega\ge\delta K^2/(2\pi)$,
so $\nu_m\ge K^2/2$, and $g(\nu_m)\ge g(K^2/2)$ since $g$ is increasing;
Theorem~\ref{thm:widesector} gives the bound.  For the limit,
$\omega\le\kappa(\delta)\lambda K/\log T$ with
$\kappa(\delta)\le\frac{2\delta}\pi(1+\log\frac\pi{2\delta})\to0$, as
$\delta\le\delta K\to0$; so $\omega=o(K)$, $\delta K^2=o(K)$,
$m\ge K-2-\omega-\delta K^2/(2\pi)=K-o(K)$, and both terms are
$(\frac12-o(1))K^2\log R$.  For the example, the coefficient of
$x^{K-k}$ has modulus $e_k(\rho_1,\ldots,\rho_K)\in[R^k,\binom Kk(2R)^k]$,
and $\sum_{k=0}^K\log\binom Kk\le(K-1)K\log2$.
\end{proof}

So for complex multiples of degree $O(K)$, as $K,R\to\infty$, both constants
are attained: $\frac12$ in $\frac12K^2\log R$ when $\delta K\to0$
(Corollary~\ref{cor:thinconst}), and $\frac\pi\delta$ in $\frac\pi\delta K\log R$ at
fixed $\delta$ (Corollary~\ref{cor:wideconst} against
Proposition~\ref{prop:widesharp}).  In between the degree bound matters:
$x^N-\rho^N$ has degree about $\pi K/\delta$, which is not $O(K)$ uniformly
in $\delta$.

\begin{corollary}[Gaussian roots in sectors of every width]\label{cor:widegauss}
In the setting of Section~\ref{sec:gaussian}, let $\varphi\in\R$,
$0<\delta<\pi/2$, $T>1$, $R\ge3T$ and $\lambda>0$, suppose that every
$\alpha_j$ has $|\alpha_j|\ge R$ and $|\arg\alpha_j-\varphi|\le\delta$, and put
$\omega=\kappa(\delta)\lambda K/\log T$.  Every nonzero $F\in\mathbb Z[x]$ divisible
by $P$ satisfies, for every integer $m\ge0$,
\[
  \Lambda(F)\ \ge\ \min\Bigl\{\lambda K\log\frac R2,\
  \frac{m(m+1)}2\log\Bigl(\frac RT-1\Bigr),\ g(\nu_m)\Bigr\},
\]
with $g$ as in Theorem~\ref{thm:widesector} and
$\nu_m=\frac\pi\delta(K-1-m-\omega)^+$ for this $\omega$.  In particular, for fixed
$\delta$, $\Lambda(F)\ge(\pi/\delta-o(1))K\log R$ as $K\to\infty$ and
$R\to\infty$.
\end{corollary}

\begin{proof}
Let $B$ be a block of any decomposition of $F$ as in
\cite[Lemma~6.2]{coefficient-mass-complex}, of degree $d$; it is
real, $|b_d|\ge1$, it vanishes at the $K$ distinct $\alpha_j$, and
$\Lambda(F)\ge\Lambda(B)$.  If $d\ge\lambda K$, \cite[Theorem~6.3]{coefficient-mass-complex}
with $\rho_{\min}\ge R>2$ gives $\Lambda(B)\ge\lambda K\log(R/2)$.  Otherwise
$\kappa(\delta)d/\log T\le\omega$, and Theorem~\ref{thm:widesector} applies to $B$;
its bound only improves when $\omega$ decreases, since $\nu_m$ then
increases and $g$ is increasing.  For the last statement take
$\lambda=\pi/\delta$ and argue as in Corollary~\ref{cor:wideconst}.
\end{proof}

The bounds of Corollary~\ref{cor:wideconst}, hence the $o(1)$ here, are
explicit in $K$ and $R$, so the limit is uniform in the $\alpha_j$.  Many
Gaussian integers cannot all have small modulus, which gives a bound
at bounded $R$.

\begin{corollary}[Many Gaussian roots]\label{cor:gausscount}
In the setting of Section~\ref{sec:gaussian}, let $\varphi\in\R$ and
$0<\delta<\pi/2$ be fixed, let $R\ge1$, and suppose that every $\alpha_j$
has $|\alpha_j|\ge R$ and $|\arg\alpha_j-\varphi|\le\delta$.  Every nonzero
$F\in\mathbb Z[x]$ divisible by $P$ satisfies
\[
  \Lambda(F)\ \ge\ \Bigl(\frac\pi\delta-o(1)\Bigr)K\log\max(R,\sqrt K)
\]
as $K\to\infty$, uniformly in $R$ and in the $\alpha_j$.  In particular
$\Lambda(F)\ge(\frac\pi{2\delta}-o(1))K\log K$ whatever $R$.
\end{corollary}

\begin{proof}
Let $K\ge3$, $X=\sqrt K/\log K$ and $R'=\max(R,X)$.  The unit squares
centred at the Gaussian integers of modulus less than $X$ are disjoint and
lie in the disc of radius $X+1$, so there are at most $\pi(X+1)^2=o(K)$ of
them.  So the set $A$ of the $j$ with $|\alpha_j|\ge R'$ has
$|A|=(1-o(1))K$.  The product of the $(x-\alpha_j)(x-\bar\alpha_j)$,
$j\in A$, divides $P$, hence $F$, and Corollary~\ref{cor:widegauss} for the
$\alpha_j$, $j\in A$, with $R'$ in place of $R$, gives
$\Lambda(F)\ge(\pi/\delta-o(1))|A|\log R'$, as $|A|\to\infty$ and
$R'\ge X\to\infty$.  Finally
$\log R'\ge\log\max(R,\sqrt K)-\log\log K$, and
$\log\log K=o(\log\max(R,\sqrt K))$.
\end{proof}

It improves on Corollary~\ref{cor:widegauss} only when $R<\sqrt K$, and
says nothing about $K^2\log R$.

\section{Integer imitations of $x^N-\rho^N$}\label{sec:imitations}

Corollary~\ref{cor:widegauss}
leaves open whether $\Lambda(F)\ge cK^2\log R$ for every integer multiple at
Gaussian roots in a wide sector.  Over the reals the bound fails because of
$x^N-\rho^N$, whose mass sits in its lowest coefficient, followed by a gap
$N$; here a polynomial $F$ with lowest nonzero coefficient $f_e$ has a
\emph{gap $m$ after its lowest coefficient} if $f_{e+1}=\cdots=f_{e+m-1}=0$,
that is, if $F=x^e(f_e+x^mG)$ for a polynomial $G$.  The two integer
imitations of $x^N-\rho^N$ are limited at Gaussian roots: a
polynomial $A((x-t)^N)$ has at most $2\deg A$ Gaussian roots in the upper
half-plane, and under the hypotheses of Proposition~\ref{prop:gausslow},
item~2, a long gap after the lowest coefficient is expensive.

\begin{lemma}[Polynomials in $(x-t)^N$]\label{lem:gaussbinom}
Let $A\in\mathbb Z[y]$ be nonzero of degree $s$, $t\in\mathbb Z$ and $N\ge1$.
Then $A\bigl((x-t)^N\bigr)$ has at most $2s$ distinct roots $\alpha\in\mathbb Z[i]$
with $\operatorname{Im}\alpha>0$.  In particular a binomial $a(x-t)^N+b$ with
$a,b\in\mathbb Z\setminus\{0\}$ has at most two, and $x^4+4$, with the roots
$\pm1+i$ in the upper half-plane, has two.
\end{lemma}

\begin{proof}
For each such root $\alpha$ the number $(\alpha-t)^N$ is one of the at most $s$
distinct roots of $A$, which sorts the roots $\alpha$ into at most $s$ classes.
Let $\alpha,\alpha'$ lie in one class, and put $\beta=\alpha-t$,
$\beta'=\alpha'-t$.  These are nonzero Gaussian integers with
$\beta'^N=\beta^N$, so $\beta'/\beta$ is a root of unity in $\mathbb Q(i)$,
that is, one of $\pm1,\pm i$.  If $\beta=u+iv$, the imaginary parts of $\beta$,
$i\beta$, $-\beta$, $-i\beta$ are $v$, $u$, $-v$, $-u$, of which at most two are
positive; since $\operatorname{Im}\alpha=\operatorname{Im}\beta$, each class has
at most two members.  Finally $x^4+4=(x^2-2x+2)(x^2+2x+2)$.
\end{proof}

So $x^N-b$, $b\in\mathbb Z$, carries at most two of the prescribed roots
$\alpha_j$, whatever $N$ and $b$, where Proposition~\ref{prop:widesharp} needs
$K$; more generally $A(x^N)$, whose mass is that of $A$
\cite[Theorem~2.2]{coefficient-mass-complex}, carries at most $2\deg A$ of them.

For the gap, recall $P(0)=\prod_j\rho_j^2$; as $P(0)\ne0$, $1/P$ is a power
series at $0$.

\begin{proposition}[The lowest coefficients]\label{prop:gausslow}
In the setting of Section~\ref{sec:gaussian}, let $m\ge1$, write
$1/P(x)=\sum_{l\ge0}u_lx^l$ as a power series at $0$, and let $d_m$ be the
least positive integer with $d_mu_l\in\mathbb Z$ for every $l<m$.
\begin{enumerate}
\item For $f\in\mathbb Z$, some $F\in\mathbb Z[x]$ divisible by $P$ has
$F\equiv f\pmod{x^m}$ if and only if $d_m\mid f$; and $d_m\mid P(0)^m$.
\item If the norms $\rho_j^2=a_j^2+c_j^2$ are odd and pairwise coprime and
$\gcd(a_j,c_j)=1$ for every $j$, then $d_m=P(0)^m$.  Consequently, if
$F\in\mathbb Z[x]$ is divisible by $P$ and $F=x^e(f_e+x^mG)$ with $f_e\ne0$
and $G\in\mathbb Z[x]$, a gap $m$ after its lowest coefficient, then $P(0)^m$
divides $f_e$, and
\[
  \Lambda(F)\ \ge\ m\log P(0)\ \ge\ 2Km\log\rho_{\min}.
\]
\item If all the norms $\rho_j^2$ equal one integer $n$, then
$d_m\mid n^{K+m-1}$.
\end{enumerate}
\end{proposition}

\begin{proof}
1.  The integers $f$ with $fu_l\in\mathbb Z$ for all $l<m$ form a subgroup of
$\mathbb Z$, namely $d_m\mathbb Z$.  Since $P$ is monic, a multiple of $P$ in
$\mathbb Z[x]$ is $F=PQ$ with $Q\in\mathbb Z[x]$, and since $P$ is a unit of
$\mathbb Q[[x]]$, $F\equiv f\pmod{x^m}$ if and only if $Q\equiv f/P\pmod{x^m}$.
If so, the $fu_l$, $l<m$, are coefficients of $Q$, hence integers, and
$d_m\mid f$.  Conversely, if $d_m\mid f$, then $Q=\sum_{l<m}fu_lx^l$ lies in
$\mathbb Z[x]$ and $F=PQ\equiv f$.  Writing $P=P(0)+xP_1$ with
$P_1\in\mathbb Z[x]$,
\[
  \frac1P=\sum_{k\ge0}\frac{(-1)^kx^kP_1^k}{P(0)^{k+1}},
\]
and only $k\le l$ contributes to $u_l$, so $P(0)^{l+1}u_l\in\mathbb Z$ and
$d_m\mid P(0)^m$.

2.  First, the $2K$ Gaussian integers $\alpha_j$, $\bar\alpha_j$ are pairwise
coprime.  A Gaussian prime $\varpi$ dividing $\alpha_j$ and $\bar\alpha_j$ divides
$\alpha_j+\bar\alpha_j=2a_j$ and $\alpha_j-\bar\alpha_j=2ic_j$, hence
$2\gcd(a_j,c_j)=2$, so $\varpi$ is an associate of $1+i$ and $2=|\varpi|^2$ divides
the odd norm $\rho_j^2$; and a common prime factor of one of $\alpha_j,\bar\alpha_j$
and one of $\alpha_l,\bar\alpha_l$, $l\ne j$, would give $|\varpi|^2>1$ dividing
both $\rho_j^2$ and $\rho_l^2$.  Now let $F\in\mathbb Z[x]$ be divisible by $P$
with $F\equiv f\pmod{x^m}$, so $F=f+x^mG$ with $G\in\mathbb Z[x]$.  For each $j$,
$0=F(\alpha_j)=f+\alpha_j^mG(\alpha_j)$ with $G(\alpha_j)\in\mathbb Z[i]$, so
$\alpha_j^m\mid f$, and conjugating, $\bar\alpha_j^m\mid f$.  By coprimality
$\prod_j\alpha_j^m\bar\alpha_j^m=P(0)^m$ divides $f$ in $\mathbb Z[i]$, hence in
$\mathbb Z$.  Taking $f=d_m$, which item~1 allows, gives $P(0)^m\mid d_m$, and
item~1 gives the reverse.  For the consequence, $P(0)\ne0$ makes $P$ divide
$F_0=f_e+x^mG$, and $F_0\equiv f_e\pmod{x^m}$, so $d_m=P(0)^m$ divides $f_e$.
Then $\Lambda(F)\ge\log|f_e|\ge m\log P(0)$, and
$P(0)=\prod_j\rho_j^2\ge\rho_{\min}^{2K}$.

3.  Let $z$ run over the $2K$ roots $\alpha_j,\bar\alpha_j$ of $P$.  Then
$\prod_z(-z)=\prod_j\alpha_j\bar\alpha_j=n^K$, so
$P(x)=n^K\prod_z(1-x/z)$ and
\[
  \frac1P=n^{-K}\prod_z\sum_{k\ge0}\frac{x^k}{z^k}
  =n^{-K}\sum_{l\ge0}h_l(z^{-1})\,x^l,
\]
where $h_l$ is the complete homogeneous symmetric polynomial of degree $l$ in
$2K$ variables.  Since $z^{-1}=\bar z/n$ and the multiset of the $\bar z$ is
that of the $z$, $h_l(z^{-1})=n^{-l}h_l(z)$.  As a symmetric polynomial with
integer coefficients in the roots of the monic $P\in\mathbb Z[x]$, $h_l(z)$ is
an integer.  So $n^{K+l}u_l\in\mathbb Z$, and $d_m\mid n^{K+m-1}$.
\end{proof}

Under the hypotheses of item~2, then, an integer multiple imitating $x^N-\rho^N$,
with a gap $m\ge\pi K/\delta$ after its lowest coefficient, pays at least
$(2\pi/\delta)K^2\log\rho_{\min}$.  The coprimality cannot be dropped: at a
common norm $n$, item~3 gives $d_m\le n^{K+m-1}$, far below
$P(0)^m=n^{Km}$ when $K,m\ge2$ (for $n=5\cdot13\cdot17\cdot29$ there are $32$
Gaussian integers of norm $n$ in the upper half-plane).  So at a common norm
some multiple has lowest $m=O(K)$ coefficients $(d_m,0,\ldots,0)$, of mass
$\log d_m\le(K+m-1)\log n$, and no bound that looks only at them gives order
$K^2\log R$ at every root set.

At Gaussian roots a small second row confines an integer multiple to one
level set of a polynomial with small nonleading coefficients, whereas over
the reals $x^N-\rho^N$ has $b_2=0$.

\begin{proposition}[One heavy coefficient]\label{prop:gaussheavy}
In the setting of Section~\ref{sec:gaussian}, assume $\rho_{\min}>2$, and let
$F\in\mathbb Z[x]$ be nonzero and divisible by $P$, with lowest nonzero
coefficient $v=f_e$.  Call a coefficient $f_i$ of $F$ \emph{heavy} if
$|f_i|\ge\rho_{\min}/2$.
\begin{enumerate}
\item For any decomposition of $F$ as in \cite[Lemma~6.2]{coefficient-mass-complex}, the
lowest coefficient of every block is heavy and nonleading, so there are at most
as many blocks as heavy nonleading coefficients of $F$.
\item If $b_2(F)<\rho_{\min}/2$, then $F=x^e(v+H)$ with $H\in\mathbb Z[x]$,
$H(0)=0$ and every nonleading coefficient of $H$ of modulus less than
$\rho_{\min}/2$; moreover $H(\alpha_j)=-v$ for every $j$, and
$(\rho_{\min}/2)^{\deg H}\le|v|$.
\end{enumerate}
\end{proposition}

\begin{proof}
1.  A block $B$ is $PQ$ with $Q\in\mathbb Z[x]$, and $B(0)\ne0$, so
$|B(0)|\ge P(0)\ge\rho_{\min}^2>\rho_{\min}/2$.  Its degree is at least $2K\ge2$,
so its lowest coefficient is not the leading coefficient of $F$, and distinct
blocks have distinct lowest positions.

2.  At most one nonleading coefficient of $F$ is heavy, so by item~1 a
decomposition of $F$ as in \cite[Lemma~6.2]{coefficient-mass-complex} has exactly one block,
$F=x^eB$, and $B(0)=v$ is the heavy one; every other
nonleading coefficient of $F$ has modulus less than $\rho_{\min}/2$.  Write
$B=v+H$; since $B(\alpha_j)=0$, $H(\alpha_j)=-v$.  Let
$d=\deg H=\deg B$.  The split $d$ of the $P$-irreducible $B$ is unclean at
some $\alpha_j$, and $\rho_j\ge\rho_{\min}>2$, so by
\cite[Lemma~6.1]{coefficient-mass-complex} some $i<d$ has
$|b_i|\ge(\rho_j/2)^{d-i}\ge\rho_{\min}/2$.  This $b_i$ is a heavy
nonleading coefficient of $F$, so $i=0$ and $|v|\ge(\rho_{\min}/2)^d$.
\end{proof}

So, when $\rho_{\min}>2$, $b_2(F)\ge\rho_{\min}/2$ unless every $\alpha_j$
lies in one level set $\{y:H(y)=-v\}$ of such an $H$.  Three Gaussian points
above the axis can share such a level set, at arbitrarily large
$\rho_{\min}$.  Let $a^2-3c^2=1$ with $a>0$ and $c\ge4$; these solutions are
$a+c\sqrt3=(2+\sqrt3)^n$, $n\ge2$, so $c=4,15,56,\ldots$.  Then
\[
  y^3-y-(8c^3+2c)i=(y+2ci)\bigl(y^2-2ciy-(4c^2+1)\bigr)
  =(y+2ci)(y-a-ci)(y+a-ci),
\]
so $y^3-y$ is $(8c^3+2c)i$ at $\pm a+ci$ and $-(8c^3+2c)i$ at $2ci$, and
$(y^3-y)^2$, with nonleading coefficients $-2$ and $1$, takes the value
$-(8c^3+2c)^2$ at the three points $\pm a+ci$, $2ci$, which have
$\rho_{\min}=2c>4$.  At $c=4$ these are $\pm7+4i$, $8i$ and the value is
$-270400$.  In the other direction, the $\alpha_j$ lie near one circle, so
there are fewer than $6\rho_{\min}$ of them.

\begin{proposition}[Near one circle]\label{prop:gausscircle}
In the setting of Proposition~\ref{prop:gaussheavy}, item~2, let
$d=\deg H$ and $Q=(3\rho_{\min}-2)/(\rho_{\min}-2)$.  Then
$\rho_{\min}\le\rho_j<\rho_{\min}Q^{1/d}$ for every $j$, and
$K<6\rho_{\min}$.  Consequently
$\Lambda(F)\ge2K\log(\rho_{\min}/2)>2K\log(K/12)$.
\end{proposition}

\begin{proof}
Write $\rho=\rho_{\min}$, $H=\sum_{i=1}^dh_iy^i$ and
$B=\max_{0<i<d}|h_i|<\rho/2$, and put $\kappa=\rho/(2(\rho-1))$, so
$\kappa<1$ as $\rho>2$, and $(1+\kappa)/(1-\kappa)=Q$.  For each $j$, with
$r=\rho_j\ge\rho$, $H(\alpha_j)=-v$ gives
\[
  |h_d\alpha_j^d+v|=\Bigl|\sum_{0<i<d}h_i\alpha_j^i\Bigr|
  \le B\,\frac{r^d-r}{r-1}<\frac{\rho}{2(\rho-1)}\,r^d\le\kappa|h_d|r^d,
\]
so $(1-\kappa)|h_d|\rho_j^d<|v|<(1+\kappa)|h_d|\rho_j^d$.  The upper bound at a
$j$ with $\rho_j=\rho$ and the lower bound at any $j$ give
$(\rho_j/\rho)^d<Q$.

Since $\rho^2$ is an integer above $4$, $\rho\ge\sqrt5$, so
$Q=3+4/(\rho-2)\le11+4\sqrt5<20$ and $\log Q<3<\frac32\rho$.  Suppose
$K\ge6\rho$.  The $2K$ distinct roots $\alpha_j,\bar\alpha_j$ of $v+H$ give
$d\ge2K\ge12\rho$, so $t=(\log Q)/d$ satisfies $t<\frac18$ and
$t<\frac3{12\rho}$.  Put $\rho'=\rho Q^{1/d}=\rho e^t$.  The unit squares
centred at the $\alpha_j$ are disjoint and lie in the upper half of the
annulus $\rho-1/\sqrt2\le|w|\le\rho'+1/\sqrt2$, so
\[
  K\ \le\ \frac\pi2\Bigl(\rho'^2-\rho^2+\sqrt2(\rho+\rho')\Bigr)
  =\frac\pi2\rho\Bigl(\rho(e^{2t}-1)+\sqrt2(1+e^t)\Bigr).
\]
Here $e^{2t}-1\le2te^{2t}<2.6\,t$ and $e^t<1.14$, so
$\rho(e^{2t}-1)<2.6\cdot\frac3{12}<0.66$ and
$K<\frac\pi2\rho(0.66+3.03)<5.8\rho$, a contradiction.  Finally
$\Lambda(F)\ge\log|v|\ge d\log(\rho/2)\ge2K\log(\rho/2)$ by
Proposition~\ref{prop:gaussheavy}, and $\rho/2>K/12$.
\end{proof}

The count uses only the annulus, whose width is of order $\rho_{\min}/d$,
and not the integrality of $H$; when $d\ge8\rho_{\min}^2$ it gives
$\rho_j^2<\rho_{\min}^2+1$, so all the $\alpha_j$ share one norm.  Whether the number
of Gaussian points above the axis in one level set of a polynomial with
small nonleading coefficients is bounded independently of $\rho_{\min}$,
and whether $\Lambda(F)\ge cK^2\log R$ holds for all integer multiples at
Gaussian roots in a wide sector, remain open.

\paragraph{Data availability.}
No datasets were generated or analyzed in this study.

\paragraph{Computer assistance.}
No proof in this paper relies on computation.  The numerical thresholds after
Theorem~\ref{thm:thinsector} describe its range and are not used elsewhere,
and the value $-270400$ after Proposition~\ref{prop:gaussheavy} is a direct
evaluation.  The file \texttt{tests/test\_sectors.py} in the
repository of~\cite{coefficient-mass} checks the statements, on seeded
random Gaussian and integer multiples and on the extremal examples, each
group with a control showing that a false variant is caught; these checks
test the statements and are not used in any proof.  Polynomial products,
divisibility, coefficient magnitudes, power-series coefficients and Sturm
counts of real zeros are computed in exact integer and rational arithmetic;
floating point enters only in evaluating logarithms, in the right-hand sides
of mass bounds and the constants $E$ and $\kappa$.

\paragraph{Use of AI tools.}
Large language model assistants (Claude, Anthropic) were used in drafting and
revising the text, the proofs and the verification code.  The author has checked the mathematics and is responsible for
the content.

\paragraph{Funding and competing interests.}
The author received no specific funding for this work and has no competing
interests to declare.

\begin{thebibliography}{99}
\small
\bibitem{coefficient-mass}
B.~Pham,
\newblock Displaced-zero certificates and the coefficient mass of polynomial
multiples,
\newblock preprint, 2026, \url{https://github.com/bangyen/coefficient-mass}.

\bibitem{coefficient-mass-complex}
B.~Pham,
\newblock Coefficient mass of polynomial multiples at complex roots,
\newblock preprint, 2026, \url{https://github.com/bangyen/coefficient-mass}.

\bibitem{erdos-turan}
P.~Erd\H{o}s and P.~Tur\'an,
\newblock On the distribution of roots of polynomials,
\newblock \emph{Annals of Mathematics} (2) 51 (1950), 105--119.
\end{thebibliography}

\end{document}
